\documentclass[12pt]{article}

\usepackage[utf8]{inputenc}
\usepackage{amsmath,amssymb,mathtools}
\usepackage{booktabs}
\usepackage{graphicx}
\usepackage{hyperref}
\usepackage[final]{microtype}
\usepackage{titling}

\textwidth=6.5in
\textheight=9in
\oddsidemargin=0in
\evensidemargin=0in
\topmargin=0in
\headheight=0in
\headsep=0in
\parskip=0.18in
\parindent=0in

\hypersetup{colorlinks=true, linkcolor=blue, citecolor=blue, urlcolor=blue}

\newcommand{\Tr}{\mathrm{Tr}}
\newcommand{\Htot}{\mathcal H_{\rm tot}}
\newcommand{\Hcore}{\mathcal H_{\rm core}}
\newcommand{\Hrad}{\mathcal H_{\rm rad}}
\newcommand{\Smicro}{S_{\rm micro}}
\newcommand{\ket}[1]{|#1\rangle}
\newcommand{\bra}[1]{\langle #1|}

\title{A Finite Super-Hagedorn Quantum Evaporator}
\author{[Author]}
\date{\today}

\begin{document}
\hypersetup{pageanchor=false}
\begin{titlingpage}
\maketitle

\begin{abstract}
We study finite non-gravitational Hamiltonian systems whose core density of
states has the Schwarzschild form $\Smicro(E)\propto E^2$ over a finite
evaporation window.  A fixed Hilbert space contains core sectors, radiation
modes, in-sector scrambling, and energy-filtered emission matrix elements.
The density-of-states ratio gives the local thermal factor for emitted quanta;
with three-dimensional radiation phase space and emission strength proportional
to the number of horizon-area degrees of freedom it also gives the Schwarzschild
power law $P\propto M^{-2}$ and lifetime $\tau\propto M_0^3$.  We then test a
matrix-free time-independent Hamiltonian in this class.  In the largest
completed run
($n=4,\ldots,7$, $q=2$, dimension $71760$), the system conserves energy to
numerical precision, the core shrinks, radiation energy grows, and the emitted
spectrum is close to the microcanonical prediction.  Removing the scrambling
term removes the observed late-time power growth in this run.  Completing the
finite-Hamiltonian result requires two checks: the explicit $p=2,\eta=1$
Schwarzschild-rate run and the radiation-chain run needed to measure Page-like
transfer of information into correlations among early and late radiation.
\end{abstract}
\end{titlingpage}
\hypersetup{pageanchor=true}
\pagenumbering{arabic}

\section{Introduction}

We take as targets the standard Schwarzschild evaporation phenomena.  For an
uncharged nonrotating black hole, the Bekenstein-Hawking entropy scales as
$S_{\rm BH}\propto M^2$ and the Hawking temperature scales as
$T_H\propto M^{-1}$ \cite{Bekenstein1973,Hawking1975}.  The heat capacity is
negative.  The outgoing radiation is locally thermal, with frequency-dependent
greybody factors and phase-space weights multiplying the thermal factor
\cite{Hawking1975,Page1976}.  The luminosity increases as the black hole loses
mass, giving the Schwarzschild scaling $P\propto M^{-2}$ and lifetime
$\tau\propto M_0^3$.  If the evaporation process is unitary, the radiation
entropy follows a Page curve and late radiation is correlated with earlier
radiation after the Page time \cite{Page1993Info,Page1993Average}.  Fast
scrambling is the expected internal dynamical scale for black holes and enters
standard information-recovery arguments \cite{HaydenPreskill2007,SekinoSusskind2008}.

Existing finite models of unitary evaporation often emphasize the
information-theoretic part of this structure.  Tensor-network models represent
information transfer from a shrinking black-hole subsystem to outgoing
radiation \cite{LeutheusserVanRaamsdonk2017}.  Random-subsystem and
random-circuit models describe closely related unitary information transfer
\cite{HaydenPreskill2007,PiroliSunderhaufQi2020}.
The model studied here keeps the explicit finite-Hilbert-space setting and
adds the thermodynamic side: energy levels, a microcanonical density of states,
emitted radiation energies, and flux generated by Hamiltonian evolution.

The question in this paper is whether these phenomena can be reproduced in an
ordinary finite quantum system once the core has the Schwarzschild density of
states,
\begin{equation}
    \rho(E)\sim \exp(cE^2)
\end{equation}
over a finite energy window.  We first derive the thermodynamic consequences
of this density of states together with the usual radiation phase-space
weight.  We then construct a fixed finite Hilbert space and a time-independent
Hamiltonian whose emission matrix elements are energy filtered.  Finally, we
compare Schwarzschild-scaling, matched linear, and no-scrambling variants, and
we formulate the radiation-chain calculation needed to measure Page-like
information flow.

The model class considered below is specified by:
\begin{itemize}
    \item a fixed finite Hilbert space containing core and radiation degrees of
    freedom;
    \item a time-independent Hamiltonian on this Hilbert space;
    \item a core spectrum with super-Hagedorn density of states,
    $\Smicro(E)\propto E^2$;
    \item radiation modes carrying definite energies;
    \item emission matrix elements that conserve energy up to a finite
    spectral width;
    \item scrambling inside each fixed core sector, with the
    no-scrambling case used as a comparison.
\end{itemize}
The term super-Hagedorn is used here in its standard comparative sense:
the density of states grows faster than $\exp(\beta_H E)$.  Related usage
appears in deformed field theories with faster-than-Hagedorn spectra
\cite{LeFlochMezei2019}.

There are three claims in the present draft.  First, the thermodynamic
consequences of $\Smicro(E)\propto E^2$ follow analytically once emission is
weighted by the radiation phase space and by an emission strength proportional
to the number of horizon-area degrees of freedom.  Second, a finite
time-independent Hamiltonian with these ingredients can
be simulated directly and already shows energy-conserving shrinkage, radiation
growth, near-thermal spectra, and late-time power growth that depends on the
scrambling term.  Third, under the standard hierarchy
$t_{\rm scr}\ll t_{\rm emit}\ll t_{\rm evap}$, the same Hamiltonian language
reduces to the shrinking-isometry picture used in Page-like information-flow
arguments; the corresponding radiation-chain calculation remains to be
completed at comparable size.

The finite-Hamiltonian result aimed at here is the following: a fixed
non-gravitational Hamiltonian, with a super-Hagedorn core spectrum, whose
unitary evolution gives the thermodynamic evaporation scalings and transfers
information from the core into correlations in the outgoing radiation.  The
status of this result is:
\begin{center}
\begin{tabular}{lp{3.8in}}
\toprule
component & status\\
\midrule
Schwarzschild thermodynamics & analytic consequence of
$\Smicro(E)\propto E^2$\\
thermal emission factor & analytic weak-coupling rate equation\\
Schwarzschild power law & analytic for $p=2$ and $\eta=1$\\
finite Hamiltonian shrinkage & completed finite-size run in Table 1\\
scrambling dependence & completed comparison with $K_{\rm scr}=0$ in Table 1\\
explicit $p=2,\eta=1$ finite run & still to be regenerated in the
time-independent Hamiltonian calculation\\
Page-like information transfer & analytic target under the scrambling hierarchy;
still to be measured in the radiation-chain calculation\\
\bottomrule
\end{tabular}
\end{center}
Generating the super-Hagedorn spectrum from simpler microscopic degrees of
freedom would be a stronger result.  It is not part of the finite-Hamiltonian
claim tested in this paper.

The paper is organized as follows.  Section 2 defines the class of
super-Hagedorn Hamiltonians.  Section 3 derives the thermodynamic
rate scalings and describes the information-theoretic measures.  Section 4
gives the time-independent Hamiltonian calculation.  Section 5 gives the
completed finite-size result.  Section 6 describes the radiation-chain
diagnostic.
Section 7 states what remains for the finite-Hamiltonian result.  Section 8 discusses the
interpretation and the relation to more microscopic approaches, including
matrix-theory black-hole models \cite{BanksFischlerKlebanovSusskind1997}.

\section{Models of Super-Hagedorn Evaporation}

We first state the model before finite-dimensional approximations are chosen.
The core is organized by an area-like variable $A$.  In continuum notation,
\begin{equation}
    \Hcore=\int^{\oplus} dA\,\mathcal H_A ,
    \qquad
    \Smicro(A)=\log \dim\mathcal H_A \simeq \gamma A .
\end{equation}
The Schwarzschild-like mass relation is
\begin{equation}
    M(A)=c_M\sqrt A ,
    \qquad
    \Smicro(M)\simeq \frac{\gamma}{c_M^2}M^2 .
\end{equation}
The symbol $A$ should be read as a coarse area label; a finite area spectrum is
introduced only after the continuum Hamiltonian has been specified.

The core Hamiltonian may be written spectrally as
\begin{equation}
    H_{\rm core}
    =
    \int dA\int dE\,
    E\,\Pi_A(E),
\end{equation}
where $\Pi_A(E)$ projects onto the microcanonical shell at area $A$ and energy
$E$, with $\Tr\Pi_A(E)=\rho_A(E)$.  The only density-of-states input needed for
the thermodynamic argument is the super-Hagedorn scaling
$\log\rho_A(E)\sim \Smicro(A)$ over the evaporation window.

The outgoing radiation is described as a field.  We take
\begin{equation}
    \Htot=\Hcore\otimes\mathcal F_{\rm out},
    \qquad
    H_{\rm tot}=H_{\rm core}+K_{\rm scr}+H_{\rm out}+H_{\rm int},
\end{equation}
with
\begin{equation}
    H_{\rm out}
    =
    \sum_s\int_0^\infty
    \omega\, b_s^\dagger(\omega)b_s(\omega)\,
    d\nu_s(\omega).
\end{equation}
Here $s$ labels radiation species, angular channels, or polarizations, and
$d\nu_s(\omega)$ is the one-particle outgoing density of states.  In
three-dimensional massless radiation the density contributes the usual
$\omega^2$ phase-space factor.

The scrambling term acts inside fixed area sectors,
\begin{equation}
    K_{\rm scr}=\int dA\,K_A,\qquad K_A:\mathcal H_A\to\mathcal H_A ,
\end{equation}
and is assumed to mix a microcanonical shell on a time scale shorter than the
emission time.

The interaction lowers the core area and creates an outgoing quantum:
\begin{align}
    H_{\rm int}
    =
    \sum_s
    \int dA\,dA'\int dE\,dE'
    \int_0^\infty d\nu_s(\omega)\,
    \Big[
    &V_s(A',E';A,E;\omega)
    \ket{A',E'}\bra{A,E}\,
    b_s^\dagger(\omega)
    \nonumber\\
    &+{\rm h.c.}\Big] .
\end{align}
The kernel is supported on $A'<A$ and is sharply peaked at energy conservation,
$E-E'=\omega$.  We parameterize the matrix elements as
\begin{equation}
    V_s(A',E';A,E;\omega)
    =
    g\,R_s(A',E';A,E)\,
    A^{\eta/2}\,
    \omega^{p/2}\,
    f_\sigma(E-E'-\omega)\,
    \chi(A',A;\omega).
\end{equation}
The factor $R_s$ mixes the available final core states, $f_\sigma$ is an
energy-conservation window approaching a delta function as $\sigma\to0$, and
$\chi$ restricts the allowed area loss.  The Schwarzschild-rate case uses
$p=2$ and $\eta=1$: $p=2$ is the three-dimensional radiation phase-space
power, while $\eta=1$ represents emission strength proportional to the
area-sized number of emitting degrees of freedom.

This continuum Hamiltonian is the target object.  The finite calculations use
a sequence of approximations:
\begin{enumerate}
    \item The area label is replaced by sectors $n=n_{\min},\ldots,n_{\max}$
    with $\dim\Hcore(n)=q^n$, $\Smicro(n)=n\log q$, and
    $M_n=\alpha\sqrt n$.
    \item The core density $\rho_A(E)$ is replaced by sampled levels
    $E_{n,a}=M_n+\epsilon_{n,a}$.
    \item The area-loss kernel is restricted to adjacent transitions
    $n\to n-1$ in the present runs.
    \item The outgoing field is either represented by a finite outgoing lead or,
    in the calculation reported below, by a compact radiation register.
    \item The scrambling operators $K_A$ are represented by fixed random
    matrices inside each finite sector.
\end{enumerate}

For the compact-register calculation,
\begin{equation}
    H_{\rm rad}^{\rm reg}=\sum_\mu \omega_\mu b_\mu^\dagger b_\mu ,
\end{equation}
and
\begin{equation}
    H_{\rm emit}^{\rm reg}
    =
    \sum_{n,a,b,\mu}
    \left[
    g^{(n,\mu)}_{ba}
    \ket{n-1,b}\bra{n,a}\,b_\mu^\dagger
    +{\rm h.c.}
    \right],
\end{equation}
with
\begin{equation}
    g^{(n,\mu)}_{ba}
    =
    g\,R^{(n,\mu)}_{ba}\,
    A_n^{\eta/2}\,
    \omega_\mu^{p/2}
    \exp\left[
    -\frac{\{\beta_n(E_{n,a}-E_{n-1,b}-\omega_\mu)\}^2}{2\sigma^2}
    \right].
\end{equation}
The register Hilbert space is truncated to at most $K_{\rm rad}$ radiation
quanta,
\begin{equation}
    \Hrad^{(K_{\rm rad})}
    =
    {\rm span}\{\ket{\{N_\mu\}}:\sum_\mu N_\mu\le K_{\rm rad}\},
\end{equation}
so the simulated operator is
\begin{equation}
    H_{\rm tot}^{(K_{\rm rad})}
    =
    P_{\le K_{\rm rad}}
    \left(H_{\rm core}+K_{\rm scr}+H_{\rm rad}^{\rm reg}
    +H_{\rm emit}^{\rm reg}\right)
    P_{\le K_{\rm rad}} .
\end{equation}
This compact-register model is a small exact-evolution proxy for the outgoing
field problem.  It tests energy conservation, shrinkage, and coarse rate
scaling.  Local thermality is expected to be more sensitive to the replacement
of an outgoing field by a finite closed register.

For comparison we also consider a finite linear core,
\begin{equation}
    M_n=\alpha_{\rm lin}n ,
\end{equation}
with $\alpha_{\rm lin}$ chosen to match the initial inverse temperature of the
Schwarzschild-scaling core.

\section{Thermodynamic and Information Measures}

The sector entropy and mass law give
\begin{equation}
    \beta_n
    =
    \frac{\log q}{\alpha(\sqrt n-\sqrt{n-1})}
    =
    \frac{\log q}{\alpha}(\sqrt n+\sqrt{n-1})
    \simeq \frac{2\log q}{\alpha^2}M_n .
\end{equation}
Thus $T_n\propto M_n^{-1}$ and the microcanonical heat capacity is negative.

Here the model starts from an emission Hamiltonian.  The coarse-grained
emission rate is obtained from the usual weak-coupling transition-rate
estimate.  Fermi's golden rule plays this role in black-hole emission
calculations that couple black-hole states to far-field modes
\cite{BrownIliesiuPeningtonUsatyuk2026}.
The random-isometry and random-circuit evaporation models cited above specify
the transfer map directly, so their central calculations concern information
flow after the map has been chosen.

\subsection{Weak-Coupling Rate Equation}

The density-of-states ratio for emitting a quantum of energy $\omega$ is
\begin{equation}
    \frac{\rho(M-\omega)}{\rho(M)}
    =
    \exp\{S(M-\omega)-S(M)\}
    \simeq e^{-\beta(M)\omega}
\end{equation}
for $\omega\ll M$.  Including the radiation phase-space weight gives the local
spectrum
\begin{equation}
    P(x)\propto x^p e^{-x},\qquad x=\beta\omega .
\end{equation}

The same result follows from the emission Hamiltonian.  For a state in sector
$n$, the weak-coupling transition rate into sector $n-1$ with one emitted
quantum of energy $\omega$ has the form
\begin{equation}
    d\Gamma_n(\omega)
    =
    2\pi
    \overline{|g^{(n,\omega)}_{ba}|^2}\,
    \rho_{n-1}(M_n-\omega)\,
    d\nu_{\rm rad}(\omega),
\end{equation}
where $d\nu_{\rm rad}(\omega)$ is the one-particle radiation density of
states.  With the normalization of $R^{(n,\mu)}_{ba}$ chosen above, the
sector-size dependence comes from the explicit area factor, the radiation
phase space, and the density-of-states ratio.  Therefore
\begin{equation}
    d\Gamma_n(\omega)
    \propto
    A_n^\eta\,\omega^p
    \frac{\rho_{n-1}(M_n-\omega)}{\rho_n(M_n)}
    d\omega
    \simeq
    A_n^\eta\,\omega^p e^{-\beta_n\omega}\,d\omega .
\end{equation}
The expected particle emission rate and emitted power are then
\begin{align}
    \Gamma_{\rm quanta}(M)
    &\propto A(M)^\eta
    \int_0^\infty d\omega\,\omega^p e^{-\beta(M)\omega}
    \propto A(M)^\eta \beta(M)^{-(p+1)},\\
    P(M)
    &\propto A(M)^\eta
    \int_0^\infty d\omega\,\omega^{p+1} e^{-\beta(M)\omega}
    \propto A(M)^\eta \beta(M)^{-(p+2)} .
\end{align}
For a Schwarzschild-scaling core, $A(M)\propto n\propto M^2$ and
$\beta(M)\propto M$.  The three-dimensional massless radiation case has
$p=2$, and the horizon-area emission strength has $\eta=1$.  Thus
\begin{equation}
    \Gamma_{\rm quanta}(M)\propto M^2 M^{-3}=M^{-1},
    \qquad
    P(M)\propto M^2 M^{-4}=M^{-2}.
\end{equation}
Since $M_n=\alpha\sqrt n$, a single downward sector step changes the mass by
\begin{equation}
    \Delta M_n=M_n-M_{n-1}\simeq \frac{\alpha^2}{2M_n}.
\end{equation}
The corresponding continuum rate equation is
\begin{equation}
    \frac{dn}{dt}\propto -\Gamma_{\rm quanta}(M_n)
    \propto -n^{-1/2},
\end{equation}
so
\begin{equation}
    n(t)^{3/2}=n_0^{3/2}-\kappa t,
    \qquad
    \tau(M_0)\propto M_0^3 .
\end{equation}
These scalings are analytic consequences of the sector entropy, the radiation
phase-space exponent, the area-strength factor, and the weak-coupling emission
Hamiltonian.

The derivation assumes a separation of scales.  Let $J_{\rm scr}(n)$ denote a
typical in-sector scrambling scale and let $g_{\rm eff}(n)$ denote the
coarse-grained emission coupling after the radiation-mode and area factors are
included.  The target regime is
\begin{equation}
    g_{\rm eff}(n)\ll J_{\rm scr}(n),
    \qquad
    t_{\rm scr}(n)\ll t_{\rm emit}(n)\ll t_{\rm evap}.
\end{equation}
The first inequality makes emission a weak perturbation of the in-sector
dynamics.  The second says that the core scrambles between emissions, while
many emissions occur over the lifetime.  In this regime the thermodynamic
calculation gives a Markovian rate equation for $n(t)$, and the
information-flow calculation can be compared to a sequence of typical
shrinking isometries.

The thermodynamic observables are
\begin{align}
    \langle n\rangle(t)&=\sum_n nP_n(t),\\
    E_{\rm rad}(t)&=\langle H_{\rm rad}\rangle,\\
    E_{\rm tot}(t)&=\langle H_{\rm tot}\rangle .
\end{align}
The positive outward current across emission edges gives the emitted flux and
power.  Binning this current in
\begin{equation}
    x=\beta\omega
\end{equation}
gives the outgoing spectrum.  The comparison distribution is the
density-of-states thermal form
\begin{equation}
    P_{\rm th}(x)\propto x^p e^{-x}.
\end{equation}
The reported spectral distance is total variation distance.

Acceleration is measured by splitting the evolution into early, middle, and
late windows and comparing the mean emitted power:
\begin{equation}
    P_{\rm late}/P_{\rm early}.
\end{equation}

For information-flow diagnostics, the initial core is purified by an inert
reference $Q$,
\begin{equation}
    \ket{\Psi(0)}
    =
    \frac{1}{\sqrt d}\sum_{a=1}^{d}
    \ket{a}_Q\ket{n_{\max},a}_{\rm core}\ket{0}_{\rm rad},
    \qquad d=q^{n_{\max}}.
\end{equation}
The radiation is represented by an outgoing record chain.  We compute second
Renyi entropies,
\begin{equation}
    S_2(A)=-\log\Tr\rho_A^2,
\end{equation}
and mutual-information combinations such as
\begin{equation}
    I_2(A:B)=S_2(A)+S_2(B)-S_2(AB).
\end{equation}
The Page-like signature is transfer of correlations from $Q$ and the remaining
core into $Q$ and the radiation, followed by correlations between early and
late radiation regions.

\subsection{Information Flow Under the Scrambling Hierarchy}

Under the hierarchy above, the core state equilibrates within $\Hcore(n)$
before the next emission event.  A single emission step is therefore
described, after coarse graining over the emission time and radiation energy,
by an isometry
\begin{equation}
    V_n:\Hcore(n)\longrightarrow
    \Hcore(n-1)\otimes \Hrad^{(1)} .
\end{equation}
The dimensions satisfy
\begin{equation}
    \dim\Hcore(n)=q^n,
    \qquad
    \dim\Hcore(n-1)=q^{n-1},
\end{equation}
so each step reduces the remaining core Hilbert-space dimension by a fixed
factor and transfers the complementary degrees of freedom to radiation.  If
the in-sector dynamics makes $V_n$ typical on the accessible microcanonical
subspace, Page and decoupling arguments imply that information about the
initial purifier $Q$ remains inaccessible in small pieces of early radiation
before the Page point and is recovered in correlations involving sufficiently
large portions of the radiation after it \cite{Page1993Info,Page1993Average,HaydenPreskill2007}.

In this limit, the radiation-chain diagnostics above have a definite target:
$I_2(Q:C)$ decreases as $\dim\Hcore(n)$ shrinks, $I_2(Q:R)$ increases as the
radiation Hilbert space grows, and early--late radiation mutual information
becomes nonzero after the radiation has enough Hilbert-space dimension to
purify the reference.  The remaining finite-Hamiltonian question is whether
the specific $H_{\rm tot}$ used here reaches this typical-isometry regime at
the simulated sizes.

\section{Time-Independent Hamiltonian Calculation}

The Hamiltonian simulation evolves
\begin{equation}
    \ket{\psi(t)}=e^{-iH_{\rm tot}t}\ket{\psi(0)}
\end{equation}
using a matrix-free Krylov representation of $H_{\rm tot}$.  The Hilbert space
is fixed throughout the evolution.

The largest completed time-independent Hamiltonian run uses
\begin{equation}
    n=4,\ldots,7,\qquad q=2,
\end{equation}
twelve radiation modes, and a radiation occupation space with at most three
quanta.  The total Hilbert-space dimension is $71760$.
This run tests the finite Hamiltonian model, the shrinking dynamics, the
microcanonical spectral comparison, and the role of scrambling.  The explicit
$p=2$, $\eta=1$ benchmark corresponding to the Schwarzschild power law is
the next thermodynamic run to regenerate with the current script.

The comparison set is:
\begin{itemize}
    \item Schwarzschild-scaling core with scrambling;
    \item linear core with matched initial temperature and scrambling;
    \item Schwarzschild-scaling core with scrambling removed.
\end{itemize}

\section{Completed Finite-Size Result}

The one-seed result is:
\begin{center}
\begin{tabular}{lccccc}
\toprule
case & $\Delta\langle n\rangle$ & $E_{\rm rad}(t_f)$
& $P_{\rm late}/P_{\rm early}$ & TV & energy drift\\
\midrule
$\sqrt n$ + scrambling & 0.734 & 2.444 & 1.289 & 0.170 & $3.7\times10^{-12}$\\
linear + scrambling    & 0.668 & 2.078 & 1.319 & 0.189 & $1.9\times10^{-12}$\\
$\sqrt n$ no scrambling & 0.788 & 2.279 & 0.956 & 0.179 & $3.6\times10^{-12}$\\
\bottomrule
\end{tabular}
\end{center}

The time-independent Hamiltonian conserves energy to numerical precision.  The core
shrinks, radiation energy grows, and the emitted spectrum stays close to the
microcanonical density-of-states prediction.  Removing scrambling leaves
energy loss in place but removes sustained late-time power growth.  The
Schwarzschild-scaling scrambled case emits the most radiation and gives the
best mean spectral distance in this run.

The linear matched comparison still accelerates at this size.  Thus the present
finite run shows that late-time power growth depends on the scrambling term.  It does
not yet isolate the super-Hagedorn entropy-energy relation as the unique cause
of acceleration at this size.

\section{Radiation Chain and Information Flow}

To measure information flow, the radiation Hilbert space is replaced by an
outgoing chain,
\begin{equation}
    \Hrad=\bigotimes_{j=1}^{K}\Hrad^{(j)} .
\end{equation}
Emission writes into the first site, and a hopping term moves radiation outward:
\begin{equation}
    H_{\rm hop}
    =
    \lambda\sum_{j,\mu}
    \left(
    \ket{\mu}_{j+1}\ket{0}_j\bra{0}_{j+1}\bra{\mu}_j
    +{\rm h.c.}
    \right).
\end{equation}
Outer sites define early radiation and inner sites define late radiation.

The desired information-flow pattern is
\begin{align}
    I_2(Q:C)&\ \hbox{decreases},\\
    I_2(Q:R)&\ \hbox{increases},\\
    I_2(E:L)&\ \hbox{increases after substantial shrinkage}.
\end{align}
These diagnostics distinguish energy loss from information transfer.  They
also test the role of scrambling in routing initial-state information into
multi-part radiation correlations.

The radiation-chain Hamiltonian is the next calculation needed to complete
the information-flow part of the evaporation phenomena.  The observables above
are fixed before the calculation: the result is positive only if the same
time-independent Hamiltonian that gives energy loss also transfers the purification
reference into multi-part radiation correlations.

\section{What Remains for the Finite-Hamiltonian Result}

Two calculations remain for the finite-Hamiltonian result.  They do not require
changing the model class.

First, the explicit Schwarzschild-rate benchmark must be run in the
time-independent Hamiltonian model.  This means using the already defined
emission rule with $p=2$ and $\eta=1$, so that the finite simulation matches
the analytic power-law derivation in Section 3.  The observables are
$\Delta\langle n\rangle$, $E_{\rm rad}(t_f)$, the binned spectrum in
$x=\beta\omega$, $P_{\rm late}/P_{\rm early}$, and energy drift.  The same run
should include the no-scrambling comparison and enough seeds or cutoffs to show
that the observed rate behavior is not a single-seed artifact.

Second, the radiation-chain calculation must be completed at comparable size.
The same core Hamiltonian and emission rule are used, but emitted quanta are
written into an outgoing chain.  The test is whether $I_2(Q:C)$ decreases,
$I_2(Q:R)$ increases, and early--late radiation correlations grow after
substantial shrinkage.  The same calculation should retain the no-scrambling
comparison, because it tests whether the information transfer depends on
internal scrambling rather than only on energy loss.

The super-Hagedorn core spectrum remains a model input throughout these
calculations.  A microscopic derivation of that spectrum would address a
different question: where the Schwarzschild density of states comes from in an
ordinary many-body system.

\section{Discussion}

The construction gives a finite super-Hagedorn evaporation model: a fixed
Hilbert space, a time-independent Hamiltonian, a core with
$\Smicro(E)\propto E^2$, and radiation coupled by energy-filtered matrix
elements.  Analytically, this entropy-energy relation gives $T\propto M^{-1}$
and negative heat capacity.  With three-dimensional radiation phase space and
emission strength proportional to the number of horizon-area degrees of
freedom it also gives the Schwarzschild power and lifetime scalings.  The
current time-independent Hamiltonian simulation shows energy-conserving
shrinkage, radiation growth, a near-thermal outgoing flux, and late-time power
growth that depends on the scrambling term.

The main conceptual input is the super-Hagedorn core spectrum.  Ordinary local
many-body systems usually have extensive thermodynamics, with entropy and
energy both scaling with the number of active degrees of freedom.  The
Schwarzschild relation requires much faster state-counting growth as a
function of energy.  Matrix theory and holography address possible microscopic
origins of black-hole-like spectra; the present construction tests the
evaporation consequences once such a spectrum is present.

The finite-Hamiltonian result would be complete after the explicit
$p=2,\eta=1$ benchmark and the radiation-chain information-flow calculation.  A
positive result would give a finite non-gravitational Hamiltonian realization
of local thermality, negative heat capacity, accelerating emission, and
Page-like information transfer.  The separate microscopic problem is to derive
the super-Hagedorn core spectrum from simpler degrees of freedom.

\appendix

\section{Numerical Implementation}

The numerical calculation follows the level of specification used in
random-circuit evaporation simulations: the Hilbert space, Hamiltonian,
truncation, random seed, and observables are fixed before time evolution
\cite{PiroliSunderhaufQi2020}.  The completed run reported in Table 1 used
\begin{equation}
    q=2,\quad n_{\min}=4,\quad n_{\max}=7,\quad
    \alpha=8,\quad K_{\max}=3,
\end{equation}
twelve radiation modes, emission coupling $g=0.05$, seed $2468$, final time
$t_f=80$, and 25 saved time points.  The radiation modes are the six
dimensionless energies
\begin{equation}
    x=\beta\omega\in \{0.5,1.0,1.5,2.0,3.0,5.0\},
\end{equation}
with two copies of each mode.  At most three radiation quanta are retained in
the finite basis.  The resulting Hilbert-space dimension is $71760$.

The three cases in Table 1 use the same radiation modes, emission coupling,
occupation cutoff, and seed.  They differ only in the core energy law and the
presence or absence of $K_{\rm scr}$:
\begin{itemize}
    \item $M_n=\alpha\sqrt n$ with scrambling;
    \item $M_n=\alpha_{\rm lin}n$ with scrambling, where $\alpha_{\rm lin}$ is
    chosen to match the initial inverse temperature of the square-root case at
    $n_{\max}$;
    \item $M_n=\alpha\sqrt n$ with $K_{\rm scr}=0$.
\end{itemize}

Time evolution is computed by a matrix-free Krylov action of $H_{\rm tot}$.
The emission and scrambling matrices are sparse random matrices drawn once at
the beginning of the run and then held fixed.  The table entries are obtained
from the following time-series observables:
\begin{center}
\begin{tabular}{ll}
\toprule
quantity & definition\\
\midrule
$\Delta\langle n\rangle$ & $\langle n\rangle(0)-\langle n\rangle(t_f)$\\
$E_{\rm rad}(t_f)$ & final expectation value of $H_{\rm rad}$\\
$P_{\rm late}/P_{\rm early}$ & ratio of mean emitted power in late and early windows\\
TV & total-variation distance to $P_{\rm th}(x)\propto x^p e^{-x}$\\
energy drift & maximum absolute drift in $\langle H_{\rm tot}\rangle$\\
\bottomrule
\end{tabular}
\end{center}

The explicit Schwarzschild-rate benchmark uses the same basis and Hamiltonian
class with $p=2$ and $\eta=1$.  The radiation-chain diagnostic specified in
Section 6 uses the same core Hamiltonian and emission rule, with the radiation
space replaced by an outgoing chain and with a reference system purifying the
initial core state.

\bibliographystyle{apsrev4-2}
\bibliography{refs}

\end{document}
